```latex
\documentclass[11pt]{amsart}

\usepackage{amsmath,amssymb,amsthm,mathrsfs}
\usepackage[margin=1.1in]{geometry}
\usepackage{hyperref}
\usepackage{enumitem}

\newtheorem{theorem}{Theorem}[section]
\newtheorem{lemma}[theorem]{Lemma}
\newtheorem{proposition}[theorem]{Proposition}
\newtheorem{conjecture}[theorem]{Conjecture}
\theoremstyle{definition}
\newtheorem{definition}[theorem]{Definition}
\newtheorem{remark}[theorem]{Remark}

\newcommand{\Z}{\mathbb Z}
\newcommand{\Q}{\mathbb Q}
\newcommand{\C}{\mathbb C}
\newcommand{\PP}{\mathbb P}
\newcommand{\NS}{\operatorname{NS}}
\newcommand{\Pic}{\operatorname{Pic}}
\newcommand{\rk}{\operatorname{rk}}
\newcommand{\op}{\oplus}
\newcommand{\perpL}{^{\perp}}

\title[A note on $H$-polarised Tyurin degenerations]{A note on $H$-polarised Tyurin degenerations and the Doran--Harder--Thompson conjecture for K3 surfaces}

\date{\today}

\begin{document}
\maketitle

\begin{abstract}
The statement requested in the prompt is, up to the specialisation $L=H$, exactly the
Doran--Harder--Thompson (DHT) conjecture for K3 surfaces as formulated by Giovenzana--Thompson.
Consequently, a complete proof is currently not available in the literature. We give a precise
reduction to \cite[Conjecture~5.10]{dfk3slpms}, record the relevant lattice-theoretic input,
and explain why the remaining step (existence of an \emph{allowable} splitting loop on the mirror
elliptic fibration with the required compatibility) is open.
\end{abstract}

\tableofcontents

\section{Lattices and K3 preliminaries}

\subsection{Lattices}
Let $H$ denote the hyperbolic plane lattice, i.e.\ the even unimodular rank $2$ lattice with Gram matrix
$\begin{pmatrix}0&1\\1&0\end{pmatrix}$.
Let $E_8$ denote the negative definite even unimodular $E_8$ root lattice.
We write the K3 lattice as
\[
\Lambda_{K3}\;\cong\; H^{\oplus 3}\oplus E_8^{\oplus 2}.
\]

\subsection{Lattice polarisations and mirrors (Dolgachev)}
We use Dolgachev's notion of $m$-admissible isotropic vectors and mirror lattices.

\begin{definition}[{\cite[Definition~3.30]{dfk3slpms} after \cite[Definition~5.4]{Dol96}}]
Let $L$ be a nondegenerate lattice and $e\in L$ a primitive isotropic element.
Let $\mathrm{div}(e)$ be the positive generator of the ideal $\langle e,L\rangle\subset \Z$.
Then $e$ is called \emph{$m$-admissible} in $L$ if $\mathrm{div}(e)=m$ and there exists an isotropic
$g\in L$ with $\mathrm{div}(g)=m$ and $\langle e,g\rangle=m$.
\end{definition}

For the definition of the mirror lattice $\check L$ attached to a primitive isotropic vector $e\in L^\perp$
in the orthogonal complement of a polarising lattice $L\subset \Lambda_{K3}$, we refer to
\cite[\S5]{Dol96} and to \cite[\S5.3]{dfk3slpms}.

\section{The mirror lattice to $H$}

The prompt fixes the mirror lattice to be $H\oplus E_8\oplus E_8$.
This is indeed the Dolgachev mirror lattice to $H$ in the standard (divisibility $1$) situation.

\begin{lemma}\label{lem:mirror-of-H}
Let $L\cong H$ be a primitive sublattice of $\Lambda_{K3}$ and let $e\in L^\perp$ be a primitive isotropic
vector which is $1$-admissible in $L^\perp$. Then the associated mirror lattice is isomorphic to
\[
\check L \;\cong\; H\oplus E_8\oplus E_8.
\]
\end{lemma}

\begin{proof}
Since $L\cong H$ is unimodular and is primitively embedded, the orthogonal complement $L^\perp\subset \Lambda_{K3}$
is also even unimodular of signature $(2,18)$, hence
\[
L^\perp \;\cong\; H^{\oplus 2}\oplus E_8^{\oplus 2}.
\]
Choose an identification $L^\perp\cong H\oplus (H\oplus E_8^{\oplus 2})$ such that $e$ lies in the first $H$.
Write the first $H$ as $\Z e \oplus \Z f$ with $\langle e,f\rangle=1$ and $e^2=f^2=0$.
Then $e^\perp\subset H\oplus (H\oplus E_8^{\oplus 2})$ equals $\Z e \oplus (H\oplus E_8^{\oplus 2})$,
and quotienting by $\Z e$ gives
\[
e^\perp/\Z e \;\cong\; H\oplus E_8^{\oplus 2}.\qedhere
\]
\end{proof}

\begin{remark}
For $L=H$ the lattice $L^\perp$ is unimodular, so every primitive isotropic vector has divisibility $1$.
Thus the case relevant to the prompt is the natural one in which $\check L\cong H\oplus E_8^{\oplus 2}$.
\end{remark}

\section{The DHT conjecture for K3 surfaces}

\subsection{Mirror pairs in the sense of Giovenzana--Thompson}
Giovenzana--Thompson define a notion of ``mirror pair'' between an $L$-polarised Tyurin degeneration
and an elliptically fibred K3 surface with an \emph{allowable} loop splitting the base; see
\cite[Definition~5.1]{dfk3slpms}, together with \cite[Definition~4.11]{dfk3slpms} for allowable loops and
\cite[Definition~4.21]{dfk3slpms} for compatibility of the lattice quasipolarisation with the splitting.

\subsection{The conjecture}
The existence statement requested in the prompt is precisely the (first direction of the) DHT conjecture
for K3 surfaces in \cite{dfk3slpms}.

\begin{conjecture}[DHT for K3 surfaces {\cite[Conjecture~5.10]{dfk3slpms}}]\label{conj:DHT}
Let $\mathcal X\to \Delta$ be an $L$-polarised Tyurin degeneration of K3 surfaces, with general fibre $X$.
Assume there is a primitive $e\in I$ which is $m$-admissible in $L^\perp$, and let $\check L=e^\perp_{L^\perp}/\Z e$
be the mirror lattice. Let $Y$ be a $\check L$-quasipolarised K3 surface. Then $Y$ admits an elliptic fibration
$\pi:Y\to \PP^1$ and there exists an allowable loop $\gamma\subset \PP^1$ such that the $\check L$-quasipolarisation
is compatible with $\pi$ and $\gamma$, and such that $\mathcal X\to \Delta$ and $\pi:Y\to \PP^1$ are a mirror pair in
the sense of \cite[Definition~5.1]{dfk3slpms}.
\end{conjecture}

The authors emphasise that the choice of the loop $\gamma$ (equivalently, the choice of birational model of the Tyurin
degeneration) is the difficult part and is not determined by the lattice data alone; see
\cite[Remark~5.11]{dfk3slpms}.

\section{Identification of the prompt with the conjecture}

We now explain formally why the prompt's requested statement is a special case of Conjecture~\ref{conj:DHT}
(and hence presently unproved).

\begin{proposition}\label{prop:prompt-equals-conj}
Let $\mathcal X\to \Delta$ be a Tyurin degeneration of K3 surfaces which is $H$-polarised in the sense of the prompt
(and hence, in particular, $H$-quasipolarised in the sense of \cite[Definition~3.19]{dfk3slpms}).
Let $Y$ be a very general $(H\oplus E_8\oplus E_8)$-quasipolarised K3 surface.

Then the conclusion demanded in the prompt (existence of an elliptic fibration $\pi:Y\to \PP^1$ and allowable loop
$\gamma$, compatible with the quasipolarisation, such that $(\mathcal X\to \Delta)$ and $(\pi:Y\to \PP^1,\gamma)$ form a
mirror pair in the sense of \cite[Definition~5.1]{dfk3slpms}) is exactly the $L=H$ case of
\cite[Conjecture~5.10]{dfk3slpms}.
\end{proposition}

\begin{proof}
Take $L=H$ in \cite[Conjecture~5.10]{dfk3slpms}.
As observed in Lemma~\ref{lem:mirror-of-H}, for any primitive $1$-admissible isotropic vector $e\in L^\perp$ the mirror
lattice is $\check L\cong H\oplus E_8\oplus E_8$.
Thus the conjecture predicts precisely that an arbitrary $\check L$-quasipolarised K3 surface $Y$ admits an elliptic
fibration and allowable splitting loop $\gamma$ compatible with the quasipolarisation, yielding a mirror pair in the
sense of \cite[Definition~5.1]{dfk3slpms}.

The additional ``very general'' hypotheses appearing in the prompt are \emph{stronger} than those assumed in
\cite[Conjecture~5.10]{dfk3slpms}; they do not change the logical content that the desired conclusion is a special case
of the conjecture.
\end{proof}

\section{Status: the requested statement is currently open}

\begin{theorem}[Current status]\label{thm:status}
As of \cite{dfk3slpms}, no proof is known of the existence of an allowable loop $\gamma$ and compatible elliptic fibration
on the mirror K3 surface producing a mirror pair in the sense of \cite[Definition~5.1]{dfk3slpms}, even under the ``very
general'' hypotheses of the prompt. In particular, the statement requested in the prompt is presently conjectural.
\end{theorem}

\begin{proof}
This is immediate from Proposition~\ref{prop:prompt-equals-conj} together with the fact that the relevant existence
statement is explicitly formulated as a conjecture in \cite[Conjecture~5.10]{dfk3slpms}. Moreover, the difficulty of the
choice of $\gamma$ is explicitly discussed in \cite[Remark~5.11]{dfk3slpms}.
\end{proof}

\begin{remark}[What \emph{can} be proved unconditionally]
The paper \cite{dfk3slpms} provides a complete \emph{definition} of mirror pairs
\cite[Definition~5.1]{dfk3slpms} and proves several compatibility results \emph{conditional on already having a mirror
pair}; see for instance \cite[Theorem~5.4]{dfk3slpms} (compatibility with the quasi del Pezzo/LG model mirror symmetry
of \cite{HT20}) and \cite[Proposition~5.3]{dfk3slpms} (a simplification of the definition in the doubly admissible case).
What remains open, and what is precisely requested in the prompt, is the \emph{existence} of the fibration and loop with
the required compatibility, i.e.\ the existence direction of the DHT conjecture.
\end{remark}

\begin{thebibliography}{ABE22}

\bibitem[ABE22]{cmek3sspt}
V.~Alexeev, A.~Brunyate, and P.~Engel,
\emph{Compactifications of moduli of elliptic {K3} surfaces: Stable pair and toroidal},
Geom.\ Topol.\ \textbf{26} (2022), no.~8, 3525--3588.

\bibitem[Dol96]{Dol96}
I.~Dolgachev,
\emph{Mirror symmetry for lattice polarized {K3} surfaces},
J.\ Math.\ Sci.\ \textbf{81} (1996), no.~3, 2599--2630.

\bibitem[GT24]{dfk3slpms}
L.~Giovenzana and A.~Thompson,
\emph{Degenerations and fibrations of {K3} surfaces: Lattice polarisations and mirror symmetry},
Preprint (arXiv:2405.12009v3), July 2024.

\bibitem[HT20]{HT20}
A.~Harder and A.~Thompson,
\emph{Pseudolattices, del Pezzo surfaces, and Lefschetz fibrations},
Selecta Math.\ (N.S.) \textbf{26} (2020), no.~1, Paper No.~9, 113 pp.

\end{thebibliography}

\end{document}
```